\documentclass[11pt]{article}
\usepackage[T1]{fontenc}
\usepackage{lmodern}
\usepackage{amsmath,amssymb,amsthm}
\usepackage[margin=1in]{geometry}
\usepackage{booktabs}
\usepackage[hidelinks,pdftitle={A Complete Hamiltonicity Atlas for Acyclic-Orientation Flip Graphs Through Six Vertices},pdfauthor={Xiaojian Zeng},pdfsubject={Finite exact classification of acyclic-orientation flip-graph Hamiltonicity},pdfkeywords={acyclic orientation, Hamilton cycle, Gray code, graph census, exact certificate}]{hyperref}
\usepackage{xurl}

\newtheorem{theorem}{Theorem}
\newtheorem{lemma}[theorem]{Lemma}
\newtheorem{proposition}[theorem]{Proposition}
\newcommand{\AO}{\mathcal{AO}}

\title{A Complete Hamiltonicity Atlas for Acyclic-Orientation Flip Graphs\\
Through Six Vertices}
\author{Xiaojian Zeng\\
PBC School of Finance, Tsinghua University, 100190\\
\texttt{xiaojianzeng@foxmail.com}}
\date{}

\begin{document}
\maketitle

\begin{abstract}
For a finite graph $G$, let $\AO(G)$ have the acyclic orientations of $G$ as
vertices, with two orientations adjacent when they differ by reversing one
edge.  We classify the Hamiltonicity of $\AO(G)$ for every connected
unlabelled loopless simple graph $G$ with $2\leq |V(G)|\leq 6$.  Under the
explicit convention that a two-item cyclic Gray code is allowed for $K_2$,
$\AO(G)$ has a Hamilton cycle if and only if the acyclic signature $\sigma(G)$
vanishes.  There are 142 graph classes: 125 signature-zero positive cases and
17 negative cases, all excluded by parity imbalance.  The result is certified
by an exhaustive canonical graph census, a complete serialization of 19,716
acyclic orientations and 52,423 flips, an explicit spanning cycle for every
positive case, and exact reconstruction verifiers.  The conclusion is only
for this finite domain and does not settle the unrestricted classification.
\end{abstract}

\section{Statement and context}

Savage, Squire, and West introduced the acyclic-orientation flip graph and
studied Hamilton cycles in it as cyclic Gray codes, obtaining positive and
negative families of examples~\cite{SavageSquireWest1993}.  M\"uhlherr and
Poullot recently formulated the general classification problem, introduced
the acyclic polynomial and its signature, and proved the parity obstruction
used below~\cite{MuehlherrPoullot2026}.  Their Appendix~D reports 22 selected
named graphs, including cases left undecided there.  It is not a census of all
connected graphs of any bounded order.  The purpose of this note is to close
the complete order-at-most-six census, with the accompanying data serving as
the case-by-case certificate.

Let $G=([n],E)$ be a labelled representative and order its edges
$e_1,\ldots,e_m$ lexicographically.  Let $D_0$ direct every edge from its
smaller to its larger endpoint.  An orientation $A$ is encoded by
$x(A)\in\{0,1\}^m$, where bit $i$ is one precisely when $e_i$ is reversed
relative to $D_0$.  Define
\[
 \operatorname{rk}_{D_0}(A)=|x(A)|,
 \qquad
 \sigma(G)=\left|\sum_{A\text{ acyclic}}
          (-1)^{\operatorname{rk}_{D_0}(A)}\right|.
\]
Changing the base orientation changes at most the sign of the sum, so its
absolute value depends only on $G$.

Throughout, a Hamilton cycle means a cyclic list containing every vertex once
and having adjacent consecutive entries, including the closing pair.  In the
two-vertex boundary case we follow the cyclic-Gray-code convention: the two
orientations of $K_2$ form an allowed two-item cyclic list.  If one instead
requires a simple cycle to have at least three vertices, $K_2$ is a trivial
signature-zero exception to the theorem.

\begin{theorem}\label{thm:atlas}
Under the preceding two-item convention, for every connected unlabelled
loopless simple graph $G$ with $2\leq |V(G)|\leq 6$,
\[
       \AO(G)\text{ has a Hamilton cycle}
       \quad\Longleftrightarrow\quad \sigma(G)=0.
\]
The complete domain contains 142 isomorphism classes, of which 125 are
Hamiltonian and 17 are non-Hamiltonian.  In particular, there is no
signature-zero counterexample in this domain.
\end{theorem}

\section{Exact census}

For each $n\in\{2,3,4,5,6\}$, the census examines all
$2^{\binom n2}$ labelled simple graphs on $[n]$, rejects the disconnected
ones, applies all $n!$ vertex permutations, and retains the lexicographically
least graph6 encoding in each orbit.  Thus neither a randomized isomorphism
test nor a pre-existing small-graph table is needed for coverage.  The
resulting class counts are shown in Table~\ref{tab:census}.  Canonical
labelling is a standard tool for reproducible graph censuses; see, for
example, McKay and Piperno~\cite{McKayPiperno2014}.

For every representative, all $2^m$ orientation masks are tested for
acyclicity by integer Kahn elimination.  As a cross-check, the resulting set
is compared with the orientations induced by all $n!$ vertex orders: a vertex
order induces an acyclic orientation, and every acyclic orientation admits a
topological order.  A pair of retained masks is a flip edge exactly when its
exclusive-or has one set bit.  The serialized atlas records the graph6 name,
edge list, every retained mask, every flip edge, and the signed and absolute
signatures.

This proof-carrying architecture is methodologically close to the canonical
graph6 census, exact witnesses, and reconstructing verifier used in a recent
finite eight-vertex graph computation~\cite{Zeng2026Whiskered}.  Here the
small bound $n\leq6$ additionally permits direct traversal of every labelled
graph and every permutation.  A second census check canonicalizes the
connected entries in an independently supplied small-graph atlas and rejects
isomorphism for every distinct same-order pair.

\begin{table}[ht]
\centering
\caption{The complete census.  The second row is exactly the Hamiltonian row;
the third row is exactly the non-Hamiltonian row.}
\label{tab:census}
\begin{tabular}{lrrrrrr}
\toprule
$|V(G)|$ & 2 & 3 & 4 & 5 & 6 & Total\\
\midrule
Connected isomorphism classes & 1 & 2 & 6 & 21 & 112 & 142\\
$\sigma(G)=0$ & 1 & 2 & 5 & 18 & 99 & 125\\
$\sigma(G)\ne0$ & 0 & 0 & 1 & 3 & 13 & 17\\
\bottomrule
\end{tabular}
\end{table}

\section{Negative cases: the parity certificate}

\begin{lemma}\label{lem:parity}
If $\AO(G)$ has a Hamilton cycle, then $\sigma(G)=0$.
\end{lemma}

\begin{proof}
Color an orientation mask by the parity of its Hamming weight.  Reversing one
edge changes that parity, so $\AO(G)$ is bipartite with even- and odd-weight
parts.  A spanning cycle alternates between the two parts and hence uses the
same number of vertices in each.  Their cardinality difference is
\[
 \sum_{A\text{ acyclic}}(-1)^{\operatorname{rk}_{D_0}(A)}.
\]
It must vanish, and therefore so must its absolute value $\sigma(G)$.
\end{proof}

The exact enumeration gives 17 nonzero-signature representatives.  Their
canonical graph6 identifiers are
\begin{center}
\small\ttfamily
C] \quad DFw \quad DLs \quad D]\{ \quad E?\string~o \quad E@vg \quad E@vo
\quad EBj? \quad EBjW\\
EBjg \quad EBz\_ \quad EBzw \quad EB\string~o \quad EJeg \quad EJnW
\quad ELvg \quad E]\string~o.
\end{center}
For each one, the verifier reconstructs the acyclic masks directly from the
edge list and obtains unequal parity classes.  Lemma~\ref{lem:parity} is thus
an exact non-Hamiltonicity certificate for every negative case; no failed
cycle search is used as evidence.

\section{Positive cases: explicit cyclic witnesses}

For every one of the remaining 125 representatives, the supplementary cycle
file contains a cyclic list of all its acyclic-orientation bit vectors.  The
lists were found by a deterministic OR-Tools CP-SAT \texttt{AddCircuit}
search (version 9.14.6206, one worker, seed zero).  Search status is not part
of the proof.  A separate verifier, which does not import OR-Tools, starts
only from each graph's order and edge list and performs the following exact
checks:
\begin{enumerate}
  \item re-enumerate all $2^m$ masks and retain exactly the acyclic ones;
  \item check that the proposed list has no repetition and its underlying set
        is exactly the reconstructed acyclic-orientation set;
  \item check that every consecutive pair, including the closing pair, has
        Hamming distance one.
\end{enumerate}
The verifier accepts 125 cycles in order counts $1,2,5,18,99$, collectively
covering 17,130 orientation vertices.  The only two-item witness is the one
for $K_2$, and its two masks differ in the unique bit.  These checks prove
Hamiltonicity for every signature-zero case independently of the optimization
search.

Combining these witnesses with Lemma~\ref{lem:parity} and the exhaustive
census proves Theorem~\ref{thm:atlas}.  Since all 17 negative cases have
nonzero signature, the set of signature-zero negative cases is empty.

\section{Reproducibility and limitations}

The supplementary material consists of the following human-readable programs
and machine-readable certificates:
\begin{itemize}
  \item \texttt{atlas-enumerate.py} and \texttt{atlas-through-6.json} generate
        and serialize the complete graph, orientation, flip, and signature
        census;
  \item \texttt{atlas-verify.py} rechecks the census and its secondary atlas
        comparison;
  \item \texttt{hamilton-search.py} and \texttt{hamilton-cycles.json} record
        witness discovery and all 125 cyclic lists;
  \item \texttt{hamilton-verify.py} reconstructs and checks every positive
        cycle and every negative parity imbalance without the search solver.
\end{itemize}
All decisive predicates use finite integers and explicit lists.  A complete
replay runs the two verification programs with Python~3; witness discovery
need not be rerun.  The result concerns only connected loopless simple graphs
on two through six vertices.  It neither classifies larger graphs nor claims
that vanishing signature is sufficient without the order bound.

\paragraph{Computational and AI assistance.}
OpenAI Codex was used to assist with literature searching, implementation and
review of the exact-enumeration scripts, witness generation, verification, and
drafting of this manuscript.  NetworkX~3.4.2 was used for a secondary atlas
cross-check, and OR-Tools CP-SAT~9.14.6206 was used only to discover positive
witnesses.  The mathematical certificate consists of the explicit finite data
and deterministic reconstruction checks, not of solver status or AI output.
The author reviewed the statement, proof, source comparisons, and generated
artifacts and is responsible for the conclusions.

\bibliographystyle{plain}
\bibliography{references}

\end{document}
